


\subsection{Influence of $\lambda_{GAN}$ and $\lambda_{PER}$ in OPP}
%%%%%%%%%%%%%%%%%%%%%%%%%%%%%%%%%%%%%%%%%%%%%%%%%%%%%%
\begin{table}[t]
\setlength{\abovecaptionskip}{0cm}
        \caption{\textbf{Influence of $\lambda_{GAN}$ and $\lambda_{PER}$ from generative learning objectives (Eq. \ref{eq:generative_objectives}).}}
    \label{tab:ablation_generative_objectives}
    
    \resizebox{\linewidth}{!}{
        \begin{tabular}{l|c|cccl}
        \rowcolor[gray]{.9}    
        \hline 
        Objectives   & Weight & $\uparrow$ PSNR & $\uparrow$ SSIM & $\downarrow$ LPIPS & $\downarrow$ FID \\ \hline \hline
        \multirow{3}{*}{$\lambda_{GAN}$} & $1e^{-4}$   &  23.49               &   0.906              &       0.097            &   42.43               \\
                                        & $1e^{-3}$  & \textbf{23.69}      & \textbf{0.910}     & \textbf{0.091} & \textbf{39.70}                  \\
                                        & $1e^{-2}$   & 23.62                 &    0.909             &     0.095               &  40.14                \\ \hline \hline
        \multirow{3}{*}{$\lambda_{PER}$} & $1e^{-3}$   &     23.64             &  0.908               &    0.105                &  51.21                \\
                                        & $1e^{-2}$  & \textbf{23.69}      & \textbf{0.910}     & \textbf{0.091} & \textbf{39.70}                  \\
                                       & $1e^{-1}$   &   23.28              &  0.902               &          0.100          &     46.70             \\ \hline
        \end{tabular}
    }
    % \vspace{-3mm}
% \vspace{-4mm}
\end{table}
%%%%%%%%%%%%%%%%%%%%%%%%%%%%%%%%%%%%%%%%%%%%%%%%%%%%%%
We study the influence of $\lambda_{GAN}$ and $\lambda_{PER}$ in Tab. \ref{tab:ablation_generative_objectives}. In line with Sec. \ref{sec:ablation_OPP}, we report the performance of the $10\%$ instances from the test set of the ShapeNet Cars category. We first fix $\lambda_{PER}$ as $1e^{-2}$ and then vary $\lambda_{GAN}$ as $\{1e^{-4}, 1e^{-3}, 1e^{-2}\}$. Obviously, $\lambda_{GAN} = 1e^{-3}$ gives the best performance. Then, we fix $\lambda_{GAN}$ as $1e^{-3}$, and vary $\lambda_{PER}$ as $\{1e^{-3}, 1e^{-2}, 1e^{-1}\}$, and find that $\lambda_{PER} = 1e^{-2}$ performs better.  We notice that the weight used in this work is much smaller than the previous GAN methods, \eg, $\lambda_{GAN}=\lambda_{PER}=5$ in \cite{Pix2NeRF_cai2022pix2nerf}. We infer that this is due to the functional difference of the GAN model between us. Specifically, in our work, the intermediate feature map $\textbf{I}_m$ already contains the coarse information provided by the OG-NeRF model, and the GAN model merely needs to refine the coarse input with learned prior. Therefore, too large weight will lead to over-imagination. While in \cite{Pix2NeRF_cai2022pix2nerf}, the captured prior by the GAN model is the main component for final outputs, therefore they require a larger weight.


\subsection{Efficiency Analysis of OPP}
%%%%%%%%%%%%%%%%%%%%%%%%%%%%%%%%


\begin{table}[t]
\setlength{\abovecaptionskip}{0cm}
    \footnotesize
     \caption{\textbf{Efficiency comparisons with PixelNeRF.} The parameter number of OPP is only slightly higher than PixelNeRF, and the additional rendering of $\hat{\textbf{I}}^{G}$ and $M$ is significantly faster than $\hat{\textbf{I}}^{N}$. We test FPS on $1$ V100 with a batch size of $1$. }
     % \vspace{-4mm}
    \label{tab:ablation_efficiency}
    \centering
    \small
    % \resizebox{\linewidth}{!}{
     \setlength\tabcolsep{6.5pt}
        \begin{tabular}{l|cccc}
        \rowcolor[gray]{.9}
        \hline
        Methods           &  Param.   & FPS ($\hat{\textbf{I}}^{N}$) & FPS ($\hat{\textbf{I}}^{G}$)   & FPS ($M$)   \\ \hline \hline
        PixelNeRF~\cite{PixelNeRF_yu2021pixelnerf}   & 15.05M  &  0.54 &   -     &   -             \\
        \textbf{OPP (ours)} & 15.46M & 0.49  &  21.50 &  11.90 \\ 
        \hline
        \end{tabular}
        % \vspace{-4mm}
\end{table}


% N + Conf.: 2.088s 
% N: 2.0039
% Conf.: 0.084s
% G: 0.0465s
%%%%%%%%%%%%%%%%%%%%%%%%%%%%%%%%
\label{Sec:efficiency}
We report the efficiency comparisons with PixelNeRF in Table \ref{tab:ablation_efficiency}. With only an additional lightweight up-sampling module and several modifications on fully connected layers, the parameter number of our OPP is only $0.41M$ higher than PixelNeRF. During the inference stage, we need to forward two times, \ie, one full-sized ($H \times W$) rendering for getting $\hat{\textbf{I}}^{N}$ (same as PixelNeRF) and confidence map $M$, and one additional $H_m \times W_m$ rendering for $\hat{\textbf{I}}^{G}$. However, the CoRF MLP and up-sampling module are rather lightweight ($0.09M$, $0.11M$), and  $\hat{\textbf{I}}^{G}$ requires substantially fewer query points than $\hat{\textbf{I}}^N$. Therefore, rendering $\hat{\textbf{I}}^{G}$ and $M$ are over $40$ and $24$ times faster than $\hat{\textbf{I}}^{N}$ ($21.50$ \textit{vs.} $0.49$ FPS and $11.90$ \textit{vs.} $0.49$ FPS). In a nutshell, the inference speed is hardly affected.


\subsection{Comparisons of Computation Cost between OPP and Recent Works}

\begin{table}[t]
\setlength{\abovecaptionskip}{0cm}
    \footnotesize
     \caption{\textbf{Comparisons of computation cost with recent works~\cite{lin2023visionnerf,gu2023nerfdiff}.} Notably, ~\cite{gu2023nerfdiff} requires additional per-scene finetuning with unknown cost. }
    \label{tab:appendix_computationwithrecent}
    \centering
    \small
     \setlength\tabcolsep{8pt}
        \begin{tabular}{l|cc}
        \rowcolor[gray]{.9}
        \hline
        Methods    &  Param.   & Traning Time  \\ \hline \hline
        VisionNeRF~\cite{lin2023visionnerf}   & 122M  &  16*A100 5Days          \\
        NeRFDiff~\cite{gu2023nerfdiff}  & 400M / 1B &  8*A100 4Days + finetune    \\
        
        \textbf{OPP (ours)} & 15M & 8*V100  3Days \\ 
        \hline
        \end{tabular}
\end{table}
Recent works~\cite{lin2023visionnerf,gu2023nerfdiff} are not listed as competitors in Tab.~\ref{tab:category_specific} since they employ much heavier architecture with more computation cost, as illustrated in Tab.~\ref{tab:appendix_computationwithrecent}. Note that, except for the larger parameter number and training cost, NeRFDiff~\cite{gu2023nerfdiff} requires additional per-scene finetuning. In contrast, built on top of PixelNeRF, our OPP is quite lightweight, and requires no per-scene finetuning.


% \section{Few-shot Extension}
% \begin{table}[t!]
     \caption{\textbf{Effectiveness of our OPP under the few-shot setting.}}
     \label{tab:fewshot_extension}
     \vspace{-4mm}
  % \resizebox{\linewidth}{!}{
% \setlength{\columnwidth}{15mm}{
\setlength\tabcolsep{9pt}
   \begin{tabular}{c|lccc}
        \rowcolor[gray]{.9}
        \hline
        \multicolumn{1}{l|}{Shot}   & Methods & $\uparrow$ PSNR            & $\downarrow$ LPIPS          & $\downarrow$ FID             \\ \hline \hline
        \multirow{2}{*}{2} & PixelNeRF~\cite{PixelNeRF_yu2021pixelnerf}     & 17.62          & 0.378          & 194.71          \\
                                & \textbf{OPP}   & \textbf{17.75} & \textbf{0.345} & \textbf{143.38} \\ \hline
        \multirow{2}{*}{3} & PixelNeRF~\cite{PixelNeRF_yu2021pixelnerf}     & 18.30          & 0.358          & 182.67          \\
                                & \textbf{OPP}   & \textbf{18.45} & \textbf{0.323} & \textbf{134.24} \\ \hline
        \end{tabular}
    % }
\end{table}
% We validate the effectiveness of our method under the few-shot setting in Table~\ref{tab:fewshot_extension} on the DTU dataset. Similar as PixelNeRF, our framework works well with few-shot settings by simply averaging inputs from difference views after the first MLP of CoRF for predicting the confidence score $\alpha$.


% \subsection{Additional Visualization of CoRF \& DPF}
% \input{Figure/Appendix_CORF}
% \input{Figure/Appendix_CORF_SRN_CHAIRS}
% % \input{Figure/Appendix_CORF_NMR}
% \input{Figure/Appendix_CORF_DTU}
% We provide additional visualization examples of CoRF and DPF in Fig. \ref{fig_appendix_corf_srn_cars}
% -
% Fig. \ref{fig_appendix_corf_dtu}.


% \subsection{Comparisons with State-of-the-art}
% % \vspace{-2mm}
% We provide additional visualization examples for ShapeNet category-agnostic setting (Fig. \ref{fig_appendix_nmr_p1_vis}-\ref{fig_appendix_nmr_p3_vis}), ShapeNet category-specific setting (Fig. \ref{fig_appendix_srn_car_vis}-\ref{fig_appendix_srn_chair_vis}), and DTU dataset (Fig. \ref{fig_appendix_dtu_vis}), respectively. 

% \input{Figure/Appendix_NMR_p1_VIS}
% \input{Figure/Appendix_NMR_p2_VIS}
% \input{Figure/Appendix_NMR_p3_VIS}

% \input{Figure/Appendix_SRN_CAR_VIS}
% \input{Figure/Appendix_SRN_CHAIR}
% \input{Figure/Appendix_DTU_VIS}


